\section{Programm: Zahlen}
\label{sec:c:zahlen}

Um einen Computer dazu zu bringen, etwas zu tun, müssen wir ihm sagen, was er tun soll. Dazu brauchen wir Variablen. Eine \textbf{Variable} ist ein Platzhalter für einen Wert. Wir können ihr einen Namen geben und ihr einen Wert zuweisen. Beispiel: Wir wollen eine Variable mit dem Namen \texttt{zahl} erstellen und ihr den Wert 42 zuweisen. In C sieht das so aus:

\cppinput{c/02_zahlen.c}

\paragraph{Erklärung:}

\begin{description}
    \item[\texttt{int zahl = 42;}]\hfill\\
        Das ist die \textbf{Deklaration}\protect\footnote{%
            Eine \textbf{Deklaration} ist eine Mitteilung an den Compiler: ``Es gibt etwas mit diesem Namen und diesem Typ''. Auch ist \codeinline{int main();} eine Deklaration, allerdings fehlt die Definition oder Implementation, also der Code.
        } (Festlegung des Namens und um was es sich handelt) und 
        \textbf{Initialisierung}\protect\footnote{%
            Die \textbf{Initialisierung} einer Variable ist die erste Zuweisung eines Wertes. Es ist generell empfehlenswert, jede Variable zu initialisieren (wenn man sich nicht sicher ist, nimmt man die \texttt{0}), da ansonsten die Variable einen zufälligen Wert haben kann (siehe dazu Kapitel~\ref{sec:c:warnungen}).
        } (die erste Zuweisung eines Wertes) einer Variablen.

        \texttt{int} ist der Datentyp, in diesem Fall eine ganze Zahl (Integer).\\
        \texttt{zahl} ist der Name der Variable.\protect\footnote{
            Variablenamen sind nicht beliebig. Es gibt eine ganze Reihe von Regeln, die man befolgen muss. Für den Anfang sollte folgendes gelten: Eine Variable beginnt mit einem Buchstaben (groß oder klein), danach folgen Buchstaben (groß und klein), Zahlen und das ``underscore'' (der Unterstrich \texttt{\_}). Wenn Ihr es genauer wissen wollt, dann lest im Anhang nach (Kap.~\ref{sec:referenz:c:namen}).
        }\\
        \texttt{42} ist der Wert, der der Variable zugewiesen wird.

        Es ist auch möglich, eine Variable ohne Initialisierung zu deklarieren:
\begin{code}[numbers=none]
  int zahl; // Deklaration ohne Initialisierung
\end{code}
        Den Wert kann man ihr dann zu einem späteren Zeitpunkt zuweisen:
\begin{code}[numbers=none]
  zahl = 42; // Zuweisung eines Wertes
\end{code} 
    \item[\texttt{printf("Die Zahl: \%d\backslash n", zahl);}]\hfill\\
        Das ist die \textbf{Ausgabe}.

        \texttt{printf} steht für \textit{print formatted} und schreibt Text
        auf den Bildschirm (genauer: in die Konsole).

        Das \texttt{\%d} ist ein Platzhalter\footnote{%
            Um genau zu sein, es ist eine \textbf{Formatierungsanweisung}. Eine ausführliche Erklärung findet Ihr im Anhang unter \manref{printf} auf S.~\manrefpage{printf}.
        } für eine ganze Zahl, die an dieser Stelle ausgegeben wird.

        Das \texttt{\backslash n} ist ein \textbf{Zeilenumbruch}.
        Es sorgt dafür, dass der Cursor nach der Ausgabe in die nächste Zeile springt.
\end{description}

\paragraph{Merke:}

\begin{itemize}[noitemsep]
    \item Eine Variable wird deklariert und initialisiert. Dabei wird ihr ein Datentyp zugewiesen.
    \item \texttt{printf()} kann auch den Inhalt von Variable in von Menschen lesbarer Form ausgeben.
\end{itemize}
